
\chapter{BNPL}\label{ch:bnpl}
Со стороны клиента BNPL (Buy Now Pay Later) выглядит как рассрочка:
клиент выбрал товар, разбил оплату на части и~сразу забрал покупку.
\highlight{В~момент оплаты провайдер BNPL в~реальном времени оценивает клиента
и~выносит кредитное решение}~--- решение об~одобрении или отказе по~заявке
на~рассрочку с~фиксированными условиями графика платежей. При одобрении
провайдер перечисляет продавцу полную сумму покупки сразу. Клиент затем
возвращает деньги провайдеру равными частями по графику.

Продавец получает выручку так же быстро, как при обычном эквайринге,
но платит за это комиссию провайдеру BNPL. В~этой модели кредитный риск
лежит на~провайдере: если клиент не~выплатит оставшиеся части, продавец
денег не~потеряет. В~альтернативной модели \enquote{BNPL на~карте}
(см.~раздел~\ref{sec:bnpl-on-card}) распределение риска иное~--- риск
несёт эмитент или участвующий кредитор, а~не сам провайдер платёжной
формы. Провайдер зарабатывает на комиссии продавца
и~иногда на пенях за просрочку, а~продавец получает рост конверсии
и увеличение среднего чека за счёт более низкого порога входа
для покупателя.

На уровне протокола шаблоны различаются тем, кто принимает
кредитное решение и~как провайдер встраивается в~поток оплаты.

\section{BNPL в~роли PSP}\label{sec:bnpl-as-psp}

В~этой модели провайдер BNPL работает как отдельный PSP (Payment Service
Provider, платёжный провайдер; см.~главу~\ref{ch:tokenization}), и~ТСП
заключает с~ним \textit{merchant agreement}~--- договор о~приёме платежей
конкретным методом, фиксирующий комиссии, обязательства по~отгрузке
и~порядок возвратов,~--- встраивает его SDK или перенаправление в~форму
оплаты и~обрабатывает обратные вызовы с~результатом кредитного решения.

Клиент на странице оплаты выбирает провайдера BNPL, и~провайдер
показывает условия: количество платежей, даты, сумму каждого платежа.
Кредитное решение опирается на~скоринг~--- алгоритмическую
оценку кредитоспособности заявителя по~внутренней модели на~основе
доступных данных. Расчёт с~продавцом идёт банковским переводом
или внутренним зачислением за~вычетом комиссии, не~карточной авторизацией;
продавец видит одну транзакцию на~полную сумму и~отгружает товар.
Когда провайдер собирает взносы с~карты клиента, эти платежи маркируются как
CIT при~первом подтверждении и~MIT installment в~последующих,
связанные между собой через NTID (Visa) или Trace ID (Mastercard).
Альтернативный путь~--- списания не~с~карты, а~напрямую с~банковского
счёта клиента: например, Affirm в~США типично собирает взносы через
ACH (Automated Clearing House, межбанковская сеть пакетных переводов;
см.~\autoref{ch:visa}), а~не картой.

\begin{figure}[tbp]
  \centering
  \includefiguresvg[width=\linewidth]{assets/figures/ch22-bnpl/bnpl-psp-purchase}
  \caption{Покупка: ТСП интегрирует BNPL как обычный платёжный
  метод~\cite{klarna-settlement-reports,tbank-dolyame-api}.}\label{fig:bnpl-psp-purchase}
\end{figure}

\begin{figure}[tbp]
  \centering
  \includefiguresvg[width=\linewidth]{assets/figures/ch22-bnpl/bnpl-psp-installments}
  \caption{Сбор взносов с~клиента: отдельная цепочка, MoR~---
  BNPL~\cite{mrc-nti-trace-id}.}\label{fig:bnpl-psp-installments}
\end{figure}

Интеграция похожа на подключение обычного PSP: API
для создания заказа, серверные уведомления (вебхуки) об изменении
статуса, API для возвратов. Спорные сценарии определяются договором
с~провайдером; правила платёжной системы здесь не применяются. Карточного
возвратного платежа между продавцом и~клиентом нет, потому что между
ними карточная сеть в~платеже не участвует.

По этой модели работают Klarna, Affirm, Tabby, а~также все
основные сервисы BNPL в~России.

\section{BNPL на карте}\label{sec:bnpl-on-card}

Второй шаблон встраивает рассрочку в~карточную инфраструктуру:
клиент расплачивается картой, а разбиение на части происходит
на стороне эмитента или отдельного кредитора внутри карточной
сети, так что продавец может вообще не знать, что покупка
оплачена в~рассрочку.

Visa и~Mastercard реализуют это по-разному. Visa Installments
строится на~наборе правил MIT (Merchant-Initiated Transaction,
транзакция по~инициативе продавца; см.~главу~\ref{ch:subscriptions}).
Первая авторизация проходит как обычная карточная операция с~хранимыми
данными (\textit{stored credential}~--- реквизиты карты, сохранённые
у~продавца или PSP с~согласия держателя для последующих списаний),
а~поле POS Environment Indicator (атрибут авторизационного сообщения,
описывающий канал и~режим инициации операции) принимает значение~I,
которое маркирует installment MIT. После неё продавец или планировщик
Visa последовательно выставляет транши по~сохранённым реквизитам.
Кредитный риск лежит на эмитенте, который решает, для каких карт
и~каких покупок доступна рассрочка. ТСП должно поддерживать
Installments API у~своего процессора, а~эмитент должен активировать
программу для своих карт. Без этого предложение рассрочки клиенту
не появится~\cite{visa-installments}.

Mastercard Installments в~основной реализации работает иначе.
Участвующий кредитор~-- банк или финтех в~роли BNPL, заключивший
договор с~Mastercard,~-- выпускает одноразовый виртуальный номер
карты (single-use virtual card number, действительный на~одну
конкретную покупку) на~отдельном BIN (Bank Identification Number,
банковский идентификационный префикс карты; см.~главу~\ref{ch:issuer-acquirer}).
Клиент расплачивается этим номером у~продавца, продавец получает
полную сумму сразу как обычную карточную транзакцию, а~кредитор
самостоятельно собирает взносы с~клиента~\cite{mc-installments}.
Продавцу не~нужно подключать отдельный API или активировать программу,
однако продавец может идентифицировать транзакцию по~выделенному BIN.

\highlight{Различие между двумя шаблонами меняет интеграционные требования.
BNPL в~роли PSP требует отдельного \textit{merchant agreement}
с~провайдером, и~добавить его как метод оплаты без нового договора
нельзя. BNPL на карте требует поддержки Installments API у~процессора
для Visa или наличия участвующего кредитора в~сети для Mastercard,
но не требует от ТСП отдельной интеграции с провайдером BNPL}.

К~2025~году поддержка флага \texttt{POS Environment Indicator~=~I}
и~маркировки installment~MIT смещается из~премиум-функций процессора
в~базовую сертификацию. Тот же сдвиг происходит у~эмитентов: вместо
запуска отдельного BNPL-продукта банки активируют installments
как опцию для~существующих карт, выбирая категории ТСП и~диапазоны
сумм. Сетевые BNPL-программы не~требуют от~продавца отдельного
\textit{merchant agreement}, и~это меняет конкурентную позицию
независимых провайдеров: для~ТСП проще включить сетевые installments
на~существующем процессоре, чем подключать новую PSP-интеграцию.

\section{Глобальный рынок BNPL}\label{sec:bnpl-global-market}

\highlight{Эпоха независимых BNPL-провайдеров без~банковской лицензии заканчивается}. Apple~Pay~Later запустился в~марте 2023~года
с~Goldman~Sachs на~бэкенде; в~июне 2024~года Apple закрыл сервис
после~15~месяцев работы~\cite{apple-pay-later-shutdown}. В~iOS~18
Apple заменил собственный BNPL встраиванием сторонних провайдеров
(Klarna, Affirm, банковские программы) в~Apple~Pay
checkout~\cite{apple-pay-installments-ios18}. Даже игрок с~миллиардным
пользовательским потоком и~устройством на~уровне ОС не~смог удержать
BNPL как монопродукт: модель работает только в~связке с~банковской
или~кредитной инфраструктурой партнёра.

Klarna прошла противоположный путь. Запущенная в~2005~году в~Стокгольме,
к~2021~году компания достигла валуации~\$45,6\,млрд, упала
до~\$6,7\,млрд в~июне 2022~года на~корреции BNPL-сектора и~10~сентября
2025~года провела IPO на~NYSE (тикер~KLAR) по~цене~\$40 за~акцию,
закрывшись в~первый день на~\$45,82 ($+15\,\%$)~\cite{klarna-ipo-2025}.
Метрики из~проспекта: 111~млн активных потребителей, 790~тыс.\ ТСП
в~26~странах, GMV~\$112\,млрд за~12~месяцев на~30~июня 2025~года.
Глобальный рынок BNPL по~прогнозам~--- около~\$23\,млрд в~2025~году
с~ожидаемым ростом до~\$90\,млрд к~2035~году, и~основной CFPB-отчёт
по~США за~декабрь 2025~года фиксирует сохранение высокой концентрации
у~четвёрки лидеров (Affirm, Klarna, Afterpay/Block,
PayPal)~\cite{cfpb-bnpl-2025}.

Победители имеют либо банковскую лицензию (Klarna банк с~2017~года
в~Швеции, Affirm работает через банков-партнёров и~получил полную
лицензию через приобретение PayBright в~Канаде), либо принадлежат
банку (PayPal, Afterpay внутри Block с~существующей расчётной
лицензией). Финтех без~собственного кредитного контура и~без
банковского зонтика расплачивается доступностью капитала и~стоимостью
комплаенса.

Региональные сегменты живут отдельно от~четвёрки лидеров.
На~Ближнем Востоке и~в~Северной Африке доминирует Tabby
(штаб-квартира в~Дубае, лицензия Saudi Central Bank, лицензия
Central Bank of the~UAE; основной поток~--- KSA, ОАЭ). В~Юго-Восточной
Азии конкурируют Atome (Сингапур, дочка финансовой группы Advance
Intelligence) и~Akulaku в~Индонезии. На~латиноамериканском рынке
выделяются Kueski в~Мексике и~Nubank через~Nu Pay в~Бразилии,
где BNPL встраивается в~существующую банковскую инфраструктуру
PIX. У~всех трёх регионов своя регуляторная траектория, и~глобальные
лидеры в~них почти не~работают: банковская лицензия и~локальная
кредитная инфраструктура остаются конкурентным барьером.

\section{BNPL в~России}\label{sec:bnpl-russia}

Рынок BNPL сформировался по модели \enquote{BNPL
в~роли PSP}. От западных аналогов его отличает то, что все
крупные сервисы принадлежат банкам или финансовым
группам, у~которых уже есть банковская лицензия и кредитная
инфраструктура. Сервисы принимают карты любых банков, включая чужие; ограничение
задаёт собственный скоринг сервиса, применяемый к~клиенту-плательщику.

Первый сервис BNPL в~России -- \enquote{Долями} от Т-Банка,
запущенный в~апреле 2021 года~\cite{tbank-dolyame}. Продавец
интегрирует его как внешний PSP, заключая отдельный договор
и~подключая API или плагин для CMS, и~налаживая отдельный поток
уведомлений~\cite{tbank-dolyame-business}. \enquote{Подели}
устроен аналогично; его формальный оператор -- ООО
\enquote{А-Четыре Технологии}, 100\,\%-дочернее общество
АО \enquote{Альфа-Банк}, а интеграция идёт через API, плагины
для CMS или платёжных
партнёров~\cite{alfa-podeli,alfa-podeli-business}. Сбер запустил
\enquote{Плати частями} как отдельный сервис BNPL, оператором
которого выступает ООО \enquote{Центр новых финансовых
сервисов}, ЦНФС, 100\,\%-дочернее общество Сбербанка. Продавцу
необходимо подключить его как самостоятельный платёжный способ
у~процессора, и~он не связан с~интеграцией
SberPay~\cite{sber-pay-parts}.

Перечисленные сервисы объединяет одна архитектура, в~которой
провайдер работает как PSP с~собственным merchant agreement
и скорингом, а~продавец получает полную сумму покупки за вычетом
комиссии, но отличаются они моделью подключения. \enquote{Долями}
и \enquote{Подели} требуют отдельного договора с~оператором,
тогда как \enquote{Сплит} от Яндекса работает внутри Yandex Pay
SDK, и~продавец подключает его через единую интеграцию с Yandex
Pay, не заключая отдельного договора BNPL~\cite{yandex-split}.
\enquote{Плати частями} подключается как отдельный метод
у~процессора.

К~2026~году в~реестр ОСР Банка России вошли операторы практически
всех перечисленных сервисов плюс операторы маркетплейс-ориентированных
рассрочек (Ozon, Wildberries) и~отдельные МФК~\cite{cbr-osr-registry-2026}.
Концентрация рынка вокруг банков и~маркетплейсов проявляется
структурно: 14~из~14~позиций реестра контролируются крупными
группами с~собственной розничной воронкой или~банковской лицензией;
независимых финтехов в~реестре нет.

\subsection*{Скоринг в~реальном времени}

\highlight{Кредитное решение по~BNPL принимается за~сотни миллисекунд
на~этапе подтверждения покупки}. У~российских сервисов скоринговая
модель опирается на~собственные данные банковской группы~--- историю
карточных операций, остатки на~счетах, ранее выданные кредиты~---
плюс реквизиты заказа (категория ТСП, сумма, доставка). Кросс-проверка
по~бюро кредитных историй происходит асинхронно, а~не на~критическом
пути: первое одобрение работает на~собственных данных эмитента-партнёра,
синхронизация с~БКИ идёт фоном.

Технически узкое место~--- интеграция с~внешним BKI-API: latency
запроса в~Объединённое кредитное бюро или НБКИ выходит за~SLA
страницы оплаты. У~банков-владельцев BNPL-сервиса это закрыто
вертикальной интеграцией: данные клиента доступны на~стороне
группы без~внешнего вызова. Независимый провайдер без~такой
вертикали либо принимает риск без~полного кредитного отчёта,
либо удлиняет страницу подтверждения до~неприемлемых для~конверсии
значений.

\section{Регулирование}\label{sec:bnpl-regulation}

До недавнего времени BNPL в~большинстве юрисдикций существовал
в~регуляторном пробеле. Формально кредитом BNPL не считается~---
процентов для клиента нет; экономически это кредитный продукт.

В~России этот пробел закрыт Федеральным законом от 31~июля
2025 года 283-ФЗ~\cite{fz-283}. Основной массив положений
вступил в~силу 1~апреля 2026~года; отдельные нормы статьи~11,
устанавливающие требования к~отчётности ОСР перед Банком России,
вступают в~силу 1~января 2027~года.
Закон вводит понятие \enquote{оператор сервиса рассрочки}, ОСР,
и устанавливает для него отдельные регуляторные требования:
регистрацию в~реестре Банка России~\cite{cbr-osr}, требования
к~капиталу, обязательное раскрытие полной стоимости рассрочки
для клиента, ограничение максимального срока и~запрет на
взимание вознаграждения с~клиента. Когда совокупная задолженность
клиента превышает порог, установленный частью~5 статьи~13 283-ФЗ
(50~тыс.\ рублей), данные передаются в~бюро кредитных историй.
Провайдер, который хочет продолжать работу после вступления
в~силу основного массива закона, должен получить статус ОСР.

На~1~апреля 2026~года Банк России включил в~реестр
14~юридических лиц~\cite{cbr-osr-registry-2026}. Из~них только
один банк~--- ПАО~\enquote{Совкомбанк}; остальные~13 позиций~---
микрофинансовые организации и~технологические компании, многие
связаны с~маркетплейсами и~экосистемами (Wildberries, Ozon,
Яндекс, Альфа-Банк, Т-Банк, Сбер). Подзаконные акты, определяющие
порядок ведения реестра, перечень документов и~состав сведений,
утверждены Указанием Банка России от~12.01.2026~№~7276-У\footnote{В~силе с~01.05.2026.}~\cite{cbr-7276u-osr}.
Структура реестра обнаруживает экосистемную природу российского
рынка BNPL: лидеры~--- не~независимые финтехи, а~подразделения
крупных банковских групп и~маркетплейсов с~собственной клиентской
воронкой.

\subsection*{Регулирование BNPL за~пределами~РФ}

С~2024 по~2026~год пробел закрывается синхронно в~нескольких
крупных юрисдикциях, и~это меняет калькуляцию провайдера BNPL,
работающего в~нескольких странах.

В~ЕС BNPL включён в~обновлённую Consumer Credit
Directive~\cite{eu-ccd-2023}; страны-члены обязаны принять
национальные акты до~20~ноября 2026~года. Обязательная оценка
кредитоспособности, единый стандарт раскрытия условий~--- Standard
European Consumer Credit Information~--- и~право клиента на~отказ
в~течение~14~дней без~причины. До~директивы BNPL в~большинстве
стран~ЕС подпадал под~исключение для~рассрочки без~процентов
из~прежней CCD~2008~года.

В~США CFPB в~интерпретирующем правиле от~22~мая 2024~года признал
BNPL-продукты с~разбиением на~четыре части кредитными картами
по~Regulation~Z, и~требования о~раскрытии условий, разрешении
споров и~возврате стали обязательными
для~провайдеров~\cite{cfpb-bnpl-interpretive-rule-2024}. Регулятор
не~стал писать отдельный режим~--- расширил существующий
до~BNPL-сценария. В~декабрьском CFPB-отчёте 2025~года отдельный
блок посвящён доле спорных операций и~эффективности возвратов
в~рамках Regulation~Z~\cite{cfpb-bnpl-2025}.

Великобритания пошла третьим путём. HM Treasury опубликовал
правительственный ответ и~план FCA-режима в~2025~году, и~FCA
получит надзорные полномочия над~BNPL по~стандартному набору
FCA-инструментов: обязательная оценка платёжеспособности,
информационные требования, режим обработки споров через FCA
и~Financial Ombudsman Service~\cite{fca-bnpl-2025}. Режим вступает
в~силу поэтапно с~2025 по~2027~год.

Австралия завершила процесс первой: Treasury Laws Amendment Act
от~декабря 2024~года расширил National Consumer Credit Protection
Act~2009 на~BNPL, признав продукты Low Cost Credit Contracts;
с~10~июня 2025~года провайдеры обязаны иметь Australian Credit
Licence и~соблюдать стандарт ответственного
кредитования~\cite{australia-bnpl-act-2025}.

\highlight{Регуляторный пробел, в~котором BNPL вырос до~глобального
рынка, закрылся за~два года в~четырёх крупных юрисдикциях}. Стоимость
комплаенса (отчётность, надзор, спор-режим, лицензирование) растёт
быстрее, чем удельный платёж клиента, и~это структурно благоприятствует
крупным провайдерам с~существующей лицензионной базой. Структура
российского реестра ОСР~--- 13~из~14~позиций под~экосистемным
или~банковским зонтиком~--- укладывается в~ту же логику.

\subsection*{Stacking и~видимость задолженности}

Сценарий, к~которому BNPL пришёл одновременно с~регуляторной
конвергенцией: stacking~--- одновременное открытие нескольких
рассрочек у~разных провайдеров на~одного клиента, при~котором
ни~один провайдер не~видит полной картины. До~передачи данных
в~бюро кредитных историй каждый принимает решение на~собственном
скоринге, и~суммарная задолженность клиента превышает уровень,
который любая отдельная модель пропустила бы.

Российское решение зашито в~ст.~13 283-ФЗ: при~совокупной
задолженности~50~тыс.\ руб.\ данные передаются в~БКИ независимо
от~провайдера. CFPB-отчёт за~2025~год отдельно фиксирует stacking
как фактор роста просрочек в~США и~рекомендует расширение
обязательной отчётности по~Regulation~Z до~полного перечня
BNPL-операций~\cite{cfpb-bnpl-2025}. Без~общей видимости рынок
повторяет траекторию субпрайм-ипотеки в~миниатюре: каждый
игрок видит свою долю риска, никто не~видит совокупную.

\subsection*{Возврат и~спор}

В~BNPL-модели~PSP спор между клиентом и~продавцом не~проходит
через карточную сеть: правил chargeback нет, и~арбитра в~лице
схемы тоже нет. Споры разрешаются по~договору между провайдером,
продавцом и~клиентом. Если клиент возвращает товар, провайдер
останавливает оставшиеся транши и~возвращает уже уплаченные;
если клиент жалуется на~качество, разбор идёт по~договорной
процедуре провайдера, а~не по~Visa/Mastercard правилам
для~диспутов (см.~гл.~\ref{ch:disputes}).

Для~ТСП это меняет операционную нагрузку: команда поддержки
работает в~интерфейсе провайдера BNPL и~в~интерфейсе карточных
диспутов параллельно; данные
не~синхронизированы. Регуляторная конвергенция 2024--2026 годов
вводит формальный режим спора~--- Regulation~Z в~США, право
на~отказ~14~дней в~ЕС, FCA-спор-режим в~Великобритании,
ответственное кредитование в~Австралии. У~ТСП появляется
обязанность поддерживать два процесса одновременно, и~стоимость
интеграции с~BNPL-провайдером растёт.

\practicenote{При выборе провайдера BNPL смотрите долю одобрений
на вашей аудитории и регуляторный статус: для России~--- наличие
в~реестре ОСР Банка России.}
